\documentclass[11pt]{article}
\usepackage[margin=1in]{geometry}
\usepackage{amsmath,amssymb,graphicx,booktabs,hyperref,xcolor}
\title{Computational Replication Report:\\ Sobral, Mandal \& Scheurer (2024),\\ \emph{Fractionalized Altermagnets} (arXiv:2410.10949)}
\author{Ollie (OpenClaw @ Argonne) --- TEXTURES-100 Wave 4}
\date{2026-07-19}
\begin{document}
\maketitle

\section{Paper summary}
The paper studies quantum-fluctuation-driven fractionalized phases near collinear altermagnetism.
It identifies a non-coplanar ``orbital altermagnet'' (altermagnetic symmetry in spin-rotation-invariant
observables), describes neighboring fractionalized spin liquids (Schwinger-boson / SU(2)$\to$Z$_2$ gauge
theory), and --- the subtitle result --- shows that the doped fractionalized altermagnet has an electronic
spectral function with \emph{split Fermi surfaces that preserve spin-rotation symmetry}: a
\emph{spin-symmetric band splitting}.

\section{Scope statement}
The full self-consistent Schwinger-boson mean field, the gauge classification of the spin liquids, and
the spinon--chargon convolution are deep many-body theory (out of scope). We implement the
Appendix-C chargon spectral model (Eqs.~C19--C21) exactly and reproduce the headline spin-symmetric
splitting plus its contrast with ordinary altermagnetism.

\section{Model (Eqs.~C19--C21)}
\begin{align*}
\epsilon_{\mathbf k}&=-t'(\cos k_x+\cos k_y)-\mu,\\
(g_{\mathbf k})_x+i(g_{\mathbf k})_y&=t\,(1+e^{ik_x}+e^{i(k_x-k_y)}+e^{-ik_y}),\\
(g_{\mathbf k})_z&=-t(\cos k_x-\cos k_y)+H_0\alpha,\qquad
E^{\pm}_{\mathbf k}=\epsilon_{\mathbf k}\pm|g_{\mathbf k}|,
\end{align*}
with $\boldsymbol g\cdot\boldsymbol\tau$ acting in \emph{sublattice} space. $t{=}1,t'{=}0.4,\mu{=}-0.3,H_0{=}0.5$.

\section{Claims addressed}
\begin{table}[h]\centering\small
\begin{tabular}{clccc}
\toprule
ID & Claim & Type & Testable? & Tested? \\
\midrule
C1 & Split Fermi surfaces with $d$-wave-anisotropic $2|g_{\mathbf k}|$ & spectrum & yes & yes \\
C2 & Spin-symmetric: each split band spin-degenerate & symmetry & yes & yes \\
C3 & Contrast: ordinary altermagnet (form factor in spin) is spin-polarized & symmetry & yes & yes \\
C4 & Self-consistent SB mean field; spin-liquid classification; spinon-chargon convolution & many-body & no & \textcolor{red}{out of scope} \\
\bottomrule
\end{tabular}
\end{table}

\section{Results vs.\ paper}
\begin{table}[h]\centering\small
\begin{tabular}{lccc}
\toprule
Quantity & Expected & This work & Verdict \\
\midrule
Band splitting $2|g_{\mathbf k}|$ & $>0$ (split FS) & $[0,\,8.06]$ & \checkmark \\
$d$-wave part of $g_z$: axis vs diagonal & max vs node & $2.0$ vs $0.0$ & \checkmark \\
Within-band spin splitting (fractionalized) & $0$ (spin-symmetric) & $2.7\times10^{-15}$ & \checkmark \\
Spin polarization, ordinary altermagnet & $\sim1$ (spin-split) & $1.000$ & \checkmark \\
\bottomrule
\end{tabular}
\end{table}

\section{Per-claim what worked / what didn't}
\textbf{C1 (worked).} $2|g_{\mathbf k}|>0$ gives split Fermi surfaces; the $d$-wave part of $g_z$
vanishes on the diagonals (node) and peaks on the axes --- altermagnetic anisotropy.
\textbf{C2 (worked --- the headline).} Because $\boldsymbol g\cdot\boldsymbol\tau$ lives in sublattice
space, every band is a spin-degenerate doublet (within-band spin splitting $\sim10^{-15}$): the
splitting preserves spin-rotation symmetry. \textbf{C3 (worked).} Placing the same $d$-wave form
factor in \emph{spin} space (ordinary altermagnet) gives spin polarization $\approx1$ --- qualitatively
different, confirming the distinction. \textbf{C4 (out of scope).} The self-consistent SB mean field,
the spin-liquid classification, and the full electron spectral function (spinon$\otimes$chargon) are
deep many-body theory.

\section{Critique of the replication}
Honest: we reproduce the paper's \emph{conceptual headline} (spin-symmetric band splitting) exactly
from the App-C chargon equations, and cleanly contrast it with an ordinary altermagnet. But this is a
BARE chargon dispersion --- the physical electron spectral function (Fig.~3) requires convolution with
the spinon Green's function, which we did not compute (Open Q1); nor did we solve the self-consistent
mean field selecting the fractionalized phase (Open Q2). C2's near-zero within-band splitting is exact
by the sublattice-vs-spin structure of the model, so it is a structural (not fitted) result. One
subtlety was handled correctly: individual degenerate eigenvectors have arbitrary $\langle S_z\rangle$,
so spin-rotation symmetry was verified via band DEGENERACY, not $\langle S_z\rangle$. LLM-judge:
PARTIAL (coverage 5, agreement 6), adopted.

\section{Verdict}
\textbf{PARTIAL} --- the headline spin-symmetric band splitting (split $d$-wave-anisotropic Fermi
surfaces with preserved spin-rotation symmetry) and its contrast with ordinary altermagnetism are
reproduced exactly from the App-C chargon model; the full Schwinger-boson mean field, spin-liquid
classification, and spinon-chargon electron spectral function are out of scope.

\section*{Open Questions}
See \texttt{report/open\_questions.json}: Q1 spinon-chargon convolution; Q2 self-consistent phase
selection; Q3 orbital-altermagnet$\to$spin-liquid melting; Q4 experimental probe of spin-symmetric
splitting; Q5 SR-invariant anomalous transport.

\end{document}
